% from ega3-1-fr.tex
\oldpage[0\textsubscript{III}]{8}
problème universel~>>. Plus généralement, nous allons voir que la notion de
\emph{limite projective} est un cas particulier de celle de foncteur
représentable. Rappelons que dans une catégorie $C$, on définit un
\emph{système projectif} par la donnée d'un ensemble préordonné $I$, d'une
famille $(A_\alpha)_{\alpha\in I}$ d'objets de $C$, et, pour tout couple
d'indices $(\alpha,\beta)$ tel que $\alpha\leq\beta$, d'un morphisme
$u_{\alpha\beta}:A_\beta\to A_\alpha$. Une \emph{limite projective} de ce
système dans $C$ est constituée par un objet $B$ de $C$ (noté
$\varprojlim A_\alpha$), et, pour chaque $\alpha\in I$, un morphisme
$u_\alpha:B\to A_\alpha$, tels que~: 1\textsuperscript{o}
$u_\alpha=u_{\alpha\beta}u_\beta$ pour $\alpha\leq\beta$~;
2\textsuperscript{o} Pour tout objet $X$ de $C$ et toute famille
$(v_\alpha)_{\alpha\in I}$ de morphismes $v_\alpha:X\to A_\alpha$, telle
que $v_\alpha=u_{\alpha\beta}v_\beta$ pour $\alpha\leq\beta$, il existe un
morphisme unique $v:X\to B$ (noté $\varprojlim v_\alpha$) tel que
$v_\alpha=u_\alpha v$ pour tout $\alpha\in I$ (T, 1.8). Ceci s'interprète de
la façon suivante~: les $u_{\alpha\beta}$ définissent canoniquement des
applications
\[
  \overline{u}_{\alpha\beta}:\operatorname{Hom}(X,A_\beta)
  \to\operatorname{Hom}(X,A_\alpha)
\]
qui définissent un \emph{système projectif} d'ensembles
$(\operatorname{Hom}(X,A_\alpha),\overline{u}_{\alpha\beta})$, et
$(v_\alpha)$ est par définition un élément de l'ensemble
$\varprojlim_\alpha\operatorname{Hom}(X,A_\alpha)$~; il est clair que
$X\mapsto\varprojlim_\alpha\operatorname{Hom}(X,A_\alpha)$ est un
\emph{foncteur contravariant} de $C$ dans $\mathrm{Ens}$, et l'existence de
la limite projective $B$ équivaut à dire que
$(v_\alpha)\mapsto\varprojlim v_\alpha$ est un \emph{isomorphisme} de
foncteurs en $X$
\[
  \varprojlim_\alpha\operatorname{Hom}(X,A_\alpha)
  \xrightarrow{\sim}\operatorname{Hom}(X,B)
  \tag{8.1.9.1}
\]
autrement dit que le foncteur
$X\mapsto\varprojlim_\alpha\operatorname{Hom}(X,A_\alpha)$ est
\emph{représentable}.
\end{env}

\begin{env}[8.1.10]
\label{0.8.1.10-fr}
\textit{Exemples. II : Objet final.} Soient $C$ une catégorie, $\{a\}$ un
ensemble réduit à un seul élément. Considérons le foncteur contravariant
$F:C\to\mathrm{Ens}$ qui, à tout objet $X$ de $C$ fait correspondre
l'ensemble $\{a\}$ et à tout morphisme $X\to X'$ dans $C$ l'unique
application $\{a\}\to\{a\}$. Dire que ce foncteur est \emph{représentable}
signifie qu'il existe un objet $e\in C$ tel que pour tout $Y\in C$,
$\operatorname{Hom}(Y,e)=h_e(Y)$ soit \emph{réduit à un élément}~; on dit
que $e$ est un \emph{objet final} de $C$, et il est clair que deux objets
finaux de $C$ sont isomorphes (ce qui permet de définir, en général à l'aide
de l'axiome de choix, un objet final de $C$ qu'on note alors $e_C$). Par
exemple, dans la catégorie $\mathrm{Ens}$, les objets finaux sont les
ensembles réduits à un élément~; dans la catégorie des \emph{algèbres
augmentées} sur un corps $K$ (où les morphismes sont les homomorphismes
d'algèbres compatibles avec les augmentations), $K$ est un objet final~;
dans la catégorie des \emph{$S$-préschémas} (I, 2.5.1), $S$ est un objet
final.
\end{env}

\begin{env}[8.1.11]
\label{0.8.1.11-fr}
Pour deux objets $X$, $Y$ d'une catégorie $C$, posons
$h'_X(Y)=\operatorname{Hom}(X,Y)$ et pour tout morphisme $u:Y\to Y'$, soit
$h'_X(u)$ l'application $v\mapsto uv$ de $\operatorname{Hom}(X,Y)$ dans
$\operatorname{Hom}(X,Y')$~; $h'_X$ est alors un \emph{foncteur covariant}
$C\to\mathrm{Ens}$, d'où l'on déduit comme dans
(\hyperref[0.8.1.2-fr]{8.1.2}) la définition d'un foncteur covariant
canonique $h':C^0\to\operatorname{Hom}(C,\mathrm{Ens})$~; un \emph{foncteur
covariant} $F$ de $C$ dans $\mathrm{Ens}$, autrement dit un objet de
$\operatorname{Hom}(C,\mathrm{Ens})$ est alors dit \emph{représentable}
s'il existe un objet $X\in C$ (nécessairement unique à isomorphisme unique
près) tel que $F$ soit \emph{isomorphe à $h'_X$}~; nous laissons au lecteur
le soin de développer les considérations <<~duales~>> des précédentes pour
cette notion, qui couvre cette fois celle de \emph{limite inductive}, et en
particulier la notion usuelle de <<~solution de problème universel~>>.
\end{env}

\subsection{Structures algébriques dans les catégories}
\label{subsection:0.8.2-fr}

\begin{env}[8.2.1]
\label{0.8.2.1-fr}
